\clearpage
\section{Unrestricted term-context filling}
\label{sec:copying-context-filling}
The structural graph-filling theorem is now proved for every context in the
original annotated grammar, with arbitrary finite variable and atom scopes.
It includes holes below binders, unary operations, and either argument of
application and binary operations. The terms of Proposition~12.2, the
Figure~16 translation and its depth choices are unchanged.

The theorem supplies one actual sole-hole address before the argument graph
is given. The translated host is binding-free. Filling this address with the
appropriately padded translation of any term recovers the original translation
of the filled term up to recursive graph isomorphism and equality of the outer
depth indices. Neither machine determinism nor execution invariance is a premise
of this structural result. Conversely, no termination equivalence follows
without the remaining semantic connection: the copying machine must respect
isomorphism and padding, and observational congruence must be established.

\subsection{Exact composition of depth padding}
\termdef{copying-cumulative-padding}
Write $P_d^n(G)$ for the original operation that raises the allowance of a graph
$G$ from $d$ to $d+n$. Equality transport changes only a type index; it does
not insert another graph constructor. Padding and one-step raising commute
with transport along $d=e$, because eliminating that equality leaves the
same expressions on both sides.




\lean{Copying/PaddingComposition}{9}{12}{PaddingComposition.pad_castDepth}

\lean{Copying/PaddingComposition}{13}{18}{PaddingComposition.raise_castDepth}

Successive padding satisfies
\[
 P_{d+n}^{m}(P_d^n(G))
   =\operatorname{cast}_{d+(n+m)=(d+n)+m}(P_d^{n+m}(G)).
\]
Induct on $m$. With no additional step the expressions coincide. At a
successor, unfold the last raising operation, apply the induction hypothesis,
and move the equality transport through raising using the preceding lemma.
The remaining transports agree because they express the same associativity
of addition. This is graph equality, not an observational claim.

Recursive graph isomorphism also survives depth equality transport by equality
elimination. Raising to a specified common allowance consists of iterated
padding followed by such a transport. Composing the previous section's padding
congruence with transport congruence therefore proves common-allowance
congruence.


\lean{Copying/PaddingComposition}{19}{26}{PaddingComposition.pad_add}

\lean{Copying/PaddingComposition}{27}{29}{PaddingComposition.castDepth_congr}

\lean{Copying/PaddingComposition}{30}{35}{PaddingComposition.padTo_congr}

An address has two independently indexed allowances: its host depth and its
local hole depth. Transporting the host depth leaves the local depth fixed.
Filling the transported address equals transporting the filled graph, and the
sole-hole certificate is preserved. Each statement follows by eliminating
the host-depth equality.

Transporting the local depth instead leaves the host graph unchanged. Filling
this transported address with a graph at the new local index equals filling
the old address with that graph transported by the inverse equality. Its
sole-hole certificate is again preserved. These definitions and proofs
separate the two indices rather than silently identifying them.


\lean{Copying/PaddingComposition}{36}{39}{PaddingComposition.Site.castDepth}

\lean{Copying/PaddingComposition}{40}{43}{PaddingComposition.plugAt_castDepth}

\lean{Copying/PaddingComposition}{44}{48}{PaddingComposition.SoleHole.castDepth}

\lean{Copying/PaddingComposition}{49}{51}{PaddingComposition.Site.castLocal}

\lean{Copying/PaddingComposition}{52}{55}{PaddingComposition.plugAt_castLocal}

\lean{Copying/PaddingComposition}{56}{60}{PaddingComposition.SoleHole.castLocal}

Suppose a site lies in a host of allowance $h$, while its local allowance is
$d+n$. Padding the host by $m$ steps yields a site with local allowance
$(d+n)+m$. Transport this local index to $d+(n+m)$; this is the extended site.
The previous padding certificate and local-transport certificate prove that
it remains sole.

Filling the extended site with $P_d^{n+m}(A)$ is recursively isomorphic to
$P_h^m$ applied to the original site filled with $P_d^n(A)$. First use the
local-transport filling equality. The additive padding equality rewrites its
replacement as $P_{d+n}^m(P_d^n(A))$. The inverse of the previous section's
iterated padding/filling isomorphism then gives the result. Thus accumulated
padding is applied to the original argument at its own allowance.


\lean{Copying/PaddingComposition}{61}{65}{PaddingComposition.extendSite}

\lean{Copying/PaddingComposition}{66}{69}{PaddingComposition.extendSite_sole}

\lean{Copying/PaddingComposition}{70}{78}{PaddingComposition.extend_filling}

For a target host allowance $N\geq h$, take $m=N-h$ in the preceding
construction, then transport the host index $h+(N-h)$ to $N$. This gives a
site in the existing common-allowance padding of the host, with local depth
$d+(n+(N-h))$. The certificate follows from extended-site and host-transport
certificate preservation. The filling law follows by transporting the
extended-site isomorphism along the same host equality. No additional padding
is charged for surrounding boxes already counted in $h$.


\lean{Copying/PaddingComposition}{79}{83}{PaddingComposition.toSite}

\lean{Copying/PaddingComposition}{84}{87}{PaddingComposition.toSite_sole}

\lean{Copying/PaddingComposition}{88}{96}{PaddingComposition.to_filling}

\subsection{Congruence of the constructor interpretation}

Recall that scoped isomorphism includes equality of the two graph depths,
followed by recursive isomorphism at that depth. For application, unpack both
argument isomorphisms and eliminate their depth equalities. The two expressions
then choose the same maximum allowance. Apply common-allowance congruence to
each argument and the original application congruence to their padded graphs.
The resulting scoped isomorphism has reflexive equality of outer depths.
For each binary operation the proof is identical with the existing binary
constructor congruence in place of application congruence.


\lean{Copying/ContextFilling}{9}{18}{ContextFilling.application}

\lean{Copying/ContextFilling}{19}{31}{ContextFilling.binary}

Every original context preserves scoped isomorphism of its argument graphs.
Induct on the context. The hole returns the given relation. Lambda and unary
constructors apply their established congruence theorems. At either argument
position of application or a binary operation, use the corresponding new
scoped congruence with the induction hypothesis in the selected position and
reflexive scoped isomorphism of the fixed sibling in the other position.
This is a structural congruence theorem; it makes no claim about execution.


\lean{Copying/ContextFilling}{32}{45}{ContextFilling.apply_iso}

\subsection{The unrestricted addressed filling theorem}
\termdef{copying-unrestricted-filling}

Fix a context $C$ whose hole has scope $(v,a)$ and whose outer scope is
$(k,h)$, an argument allowance $d$, and a hole name. Let $T_C$ denote the
original translated holed graph, and let $D_C$ denote its outer allowance.
There exist a natural number $n$ and a sole site $s$ of $T_C$, with local
allowance $d+n$ and interface
\[
 [\mathsf{term}]\longrightarrow
 \mathsf{environmentInterface}(v,a),
\]
such that, for every fitting graph $A$ at allowance $d$,
\[
 \bigl(D_C,\operatorname{plugAt}(s,P_d^n(A))\bigr)
 \quad\mathrel{\mathsf{ScopedIso}}\quad
 \operatorname{apply}(C,(d,A)).
\]
The quantifier over $A$ follows the choices of $n$ and $s$: both are fixed by
the context, allowance and hole name. In particular, the proof cannot choose
a different occurrence after inspecting the argument graph.

\paragraph{Hole.} Choose $n=0$ and the hole corolla's unique edge. The existing
corolla replacement isomorphism recovers $A$, and the singleton carrier proves
the sole-hole certificate.

\paragraph{Lambda.} Retain the induction hypothesis's $n$. Lift its address
through the original abstraction clause. Its sole-hole law supplies the
certificate. Compose the abstraction filling isomorphism with abstraction
congruence applied to the induction hypothesis. The host allowance increases
by one; the hole's local allowance does not.

\paragraph{Unary operation.} Retain $n$, lift the address through the original
unary clause, and use its sole-hole law. Compose unary filling with unary
congruence of the induction hypothesis. The operation and distributor are
unchanged.

\paragraph{Application, left argument.} Let the selected subcontext host have
allowance $D$ and the fixed sibling term have depth $q$. Set $N=\max(D,q)$.
Pad the fixed sibling to $N$; it is hole-free because term translation and
padding preserve this property. Lift the induction hypothesis's site to
allowance $N$ using the new common-allowance site. Its local padding count is
now $n+(N-D)$. Lift that address into the application's left argument. The
padding and application certificates prove that it remains sole.

For a graph $A$, first use application filling, then application congruence
with the common-allowance filling isomorphism in the left position and the
identity in the fixed position. This yields application of the padded original
subcontext filling and the fixed sibling. Finally use scoped application
congruence with the induction hypothesis to reach the original constructor
interpretation. The intermediate scoped isomorphism ensures that the maximum
depth chosen on the target side is exactly the same $N$.

\paragraph{Application, right argument.} Choose $N=\max(q,D)$ and the same
deficit $N-D$ for the selected subcontext. Use the right-site and right-filling
laws. Place the identity and induction isomorphisms in the opposite positions.
The ordered pair boundary is still the original one.

\paragraph{Binary operation, either argument.} Repeat the corresponding left
or right application argument with that binary operation's filling and
congruence laws. The operation parameter remains unrestricted. These seven
cases exhaust the original context grammar; no restricted-context predicate
is required.


\lean{Copying/ContextFilling}{46}{129}{ContextFilling.context_filling}

For a term $t$, instantiate the theorem at $d=\mathsf{depth}(t)$ and take
$A$ to be its unchanged translation. The previously proved constructor
interpretation identity identifies the target with the original translation
of $C[t]$, including its exact outer depth. The existing theorem that every
translated term context is binding-free supplies that certificate. Thus the
same site is sole, its host is binding-free, and its actual filling has the
required structural correspondence for arbitrary scopes and argument terms.

The graph-context interface can also be packaged using this very site, rather
than selecting a potentially different site merely from hole existence.
Combine the translated host, the constructed site, its sole-hole proof and
the host's binding-free proof into a graph context. Its graph is literally
equal to the original translated host, and the universal filling law is the
one just proved. This packages a checked structural interface; it introduces
no assumption about a copying execution.


\lean{Copying/ContextFilling}{130}{143}{ContextFilling.context_term_filling}

\lean{Copying/ContextFilling}{144}{152}{ContextFilling.context_graph_filling}

\subsection{Mixed-context regression and semantic boundary}

The fixed sibling term is $\lambda x.\lambda y.y$, represented with the
original scoped variable indices, in arbitrary ambient environments. The
mixed context has an outer lambda and dereference, and its path to the hole
passes through left application, right binary addition, right application,
and left binary assignment. Its fixed siblings include the two-box term and
unit. Hence every context constructor occurs on one path, and the hole's
scope includes the variable bound by the outer lambda.

For a hole allowance of zero the translated host has depth three: the
binary/application portion meets a depth-two sibling, and the outer lambda
adds one. This equality is definitional, so the regression genuinely includes
a maximum-depth mismatch. A second definitional equality checks that term
filling retains the exact displayed constructor expression. The unrestricted
term-filling theorem then proves the actual sole-hole, binding-free graph
correspondence for every argument term and every ambient variable/atom scope.
This is a structural regression, not a claim that the mixed expression
terminates or is well behaved under a store machine.

Appendix~B.1 item~3 no longer contains unrestricted addressed correspondence
as an unfinished task. It still contains observational term-context congruence
and the observation/padding audit. In particular, arguments of different
minimal depths must be compared through common representations, and the
concrete execution must discharge its isomorphism and padding invariance
obligations. The copying machine and the template proofs remain pending;
Proposition~12.2 is not yet proved.


\lean{Copying/ContextFillingExamples}{8}{11}{ContextFillingExamples.twoBoxSibling}

\lean{Copying/ContextFillingExamples}{12}{21}{ContextFillingExamples.mixedArgumentContext}

\lean{Copying/ContextFillingExamples}{22}{25}{ContextFillingExamples.mixedArgumentContext_depth}

\lean{Copying/ContextFillingExamples}{26}{34}{ContextFillingExamples.mixedArgumentContext_fill}

\lean{Copying/ContextFillingExamples}{35}{44}{ContextFillingExamples.mixedArgumentContext_graph_filling}
